\documentclass[10pt,twocolumn]{icfd2026paper}
\usepackage[dvipdfmx]{graphicx}
\usepackage{ascmac}
\usepackage{amsmath,amssymb}
\usepackage{indentfirst}
\usepackage{newtxtext,newtxmath} %change font 2025/05/09

\newcommand{\bm}[1]{\mbox{\boldmath{${#1}$}}}
\newcommand{\pd}[2]{{\frac{\partial {#1}}{\partial {#2}}}}
\newcommand{\pdd}[2]{{\frac{\partial^2 {#1}}{\partial {#2}^2}}}
\newcommand{\od}[2]{{\frac{{\rm{d}} {#1}}{{\rm{d}} {#2}}}}
\newcommand{\odd}[2]{{\frac{{\rm{d}}^2 {#1}}{{\rm{d}} {#2}^2}}}


%%% Title %%% 
% Please capitalize the first letter of each major word in the title.
\title{Importance of N2Hx Reaction Pathways in Modelling of NH3/H2/Air Flame Chemistry}

%%% Names of author(s): Firstame Familyame %%%
% Please underline the presenter using \underline
\author{\underline{Aahil Ashraf Urothiyil}${}^1$,
	Abe Shunnosuke${}^2$, Matsuda Dai${}^2$, Kitagawa Toshiaki${}^2$, Ekenechukwu Chijioke Okafor${}^2$
	%%% Email address of the corresponding author %%%
	%\footnote for corr. author & e-mail should be placed after the last author
	%
	\footnote{\normalsize Corresponding author: Aahil Urothiyil}
	\footnote{\normalsize {\textit{E-mail address}}:  urothiyil.aahil.412@s.kyushu-u.ac.jp}
	\\
	%%% Affiliation(s) and address %%%
	${}^1$Department of Mechanical Engineering, Kyushu University \\
	7-4-4, Motooka, Nishi-ku, Fukuoka, 819--0395, Japan \\
	% 2-1-1, Katahira, Aoba-ku, Sendai, Miyagi 980--8577, Japan \\
	${}^1$Student \\ % Affiliation
	Street, City, State, Zip-code, Country
}
%%%
\abstract{

	$(*1)$Due to the growing interest in ammonia blends as a combustion fuel, Okafor et al.\cite developed a reduced
	reaction mechanism that accurately modelled the laminar burning velocity (LBV) of the CH4/NH3/Air flames.
	$(*2)$However, for NH3/H2/Air flames, this reaction mechanism underpredicts the LBV at higher hydrogen
	concentrations. This discrepancy stems from the lack of N$_2$H$_x$ reaction pathways in the mechanism active during
	NH3/H2 combustion. $(*3)$An improved reaction was developed including the important elementary reactions and was
	validated with experimental data from existing literature.

	$(*1)$To study the combustion of ammonia blends in various applications, Okafor et al. developed a reduced reaction
	mechanism that accurately modelled the LBV of the CH4/NH3/Air flames.

	$(*2)$This reaction mechanism was investigated for how it models LBV of NH3/H2/Air flame and was modified to produce
	more accurate predictions, inline with experimental data sourced from existing literature.

	$(*2)$On the other hand, LBV predictions for NH3/H2/Air flames at higher hydrogen concentrations with this mechanism
	were not representative of actual flames. This stems from the reaction mechanism lacking N$_2$H$_x$ reaction
	pathways.

	$(*3)$An improved reaction mechanism was developed which includes N$_2$H$_x$ reaction pathways, then was validated
	with experimental data from previous studies.

}
\pagestyle{empty}

\begin{document}

\maketitle
\thispagestyle{empty}

\section{Introduction}

% Ammonia has high volumetric density, storage easier than hydrogen \cite{El-Adawy24}
% Ammonia has low LHV, flamability limit and LBV than H2 and CH4 \cite{Alnajideen24}
% Ammonia has existing infrastructure to produce, store and transport.\cite{Valera-Medina21}

Replacing fossil fuels as the primary source of energy is being increasingly more important as the environmental
impact from carbon emissions, such as global warming, continue to intensify. Ammonia and ammonia-fuel blends are
potential alternative fuels to replace traditional fossil fuels. Ammonia is an excellent hydrogen carrier, another
alternative fuel, and the combustion of ammonia produces no carbon products. Moreover, the production, storage and
transportation of ammonia is more practical relative to other carbon free fuels such as hydrogen. This along with the
existing infrastructure for producing and transporting ammonia are strong reasons for promoting ammonia's use as a fuel.

However, ammonia combustion poses certain challenges in various applications --- ammonia has a slower LBV, smaller
lower heating value and narrower flamability limit than fuels such as methane and hydrogen. These demerits need to be
overcome for ammonia to be functionally applicable in combustion systems. With regards to emissions, although no carbon
dioxide is produced from ammonia combustion, fuel NOx production becomes significant. Understanding NOx production
pathways is required to control these emissions.

Improving the combustion properties of ammonia can be achieved by mixing more combustible fuels, like methane and
hydrogen, to ammonia. Ideally, such fuel blends should be able to exhibit combustion properties similar to that of
traditional fuels while significantly reducing carbon emissions from combustion. Adding fuels with higher LBV than
ammonia's, like methane and hydrogen, can improve the LBV and other properties of the flame. Meanwhile, carbon emissions
can be kept to a minimum. Therefore, finding the exact ratios of fuels in the ammonia blends at which these flames
properties align with system requirement is neccessary and is possible with numerical simulations with reaction
mechanisms.

The aim of this study is to develop a reaction mechanism to accurately model flames for ammonia and ammonia blends,
namely NH3/H2/Air and CH4/NH3/Air flames. This mechanism prioritizes LBV, a critical flame characteristic that dictates
fuel compatability for a wide range of practical applications. Most existing reaction mechanisms from literature are
only able to reproduce some fuels' LBV accurately under select pressure, initial temperature and equivalence ratios.
Okafor reduced mechanism, for example, is designed to predict CH4/NH3/Air's LBV. This mechanism also exhibits high
reproducibility for LBVs of NH3/Air flames and H2/Air flames however underpredicts LBV of NH3/H2/Air flames. In this
study, the reason for the underprediction was studied and a mechanism was developed around the findings.

\section{Developing Reaction Mechanism}

The new reaction mechanism developed in this paper is based on the Okafor reduced mechanism. To study how ammonia and
hydrogen interact and affect LBV, Otomo mech, Nakamura mech and Gotama mech were studied. The figure below [Figure lbv
nh3/h2] comparing the different numerical simultions with the mentioned reaction mechanisms to experimental data. For
NH3/H2/Air flames, LBV predictions by Okafor reduced mech are accurate for smaller hydrogen concentrations but
underpredict otherwise. On the other hand other mechanism are able to reproduce this flame property across various
hydrogen concentrations.

\subsection{Methodology}

Numerical simulations for this study is conducted using ANSYS Chemkin-PRO 2019 R3 and ANSYS Chemkin 2025 R2. The
"Premixed Laminar Flame Speed Calculation" model was used and the initial conditions kept constant for all the
simulations were the pressure at 1.0 atm, temperature at 298.0 K, and equivalence ratio at 1.0. GRAD and CURV parameters
were set to 0.01, the computational domain was 10 cm, multicomponent transport was enabled and thermal diffusion was
enabled.

From the numerical simulations, integrated rate of production (IROP) and sensitivity of each elementary reaction was
evaluated. With the IROP, reaction pathways can be created for each mechanism. Utilizing the sensitivity analyses and
reaction pathways, the reactions impacting LBV are highlighted by recognising the differences between Okafor reduced mech
and the other mechanisms.

\subsection{Findings}

[Figures reaction pathways]

Comparing Okafor reduced mech to Nakamura mech, few key pathways depict LBV prediction disparity between the mechanisms.
NH$\rightarrow$ N$\rightarrow$ N2 and  reaction pathway has higher IROP in Nakamura mech

\section{Validation}

\section{Concluding Remarks}
The due date of the paper submission is \textbf{July 21, 2026}.

\begin{thebibliography}{99}

\bibitem{Okafor19}
E. C. Okafor, Y. Naito, S. Colson, A. Ichikawa, T. Kudo, A. Hayakawa, and H. Kobayashi,
\textit{Combustion and Flame}, 204, (2019), 162–175.

\bibitem{Valera-Medina21}
A. Valera-Medina, F. Amer-Hatem, A. K. Azad, I. C. Dedoussi, M. De Joannon, R. X. Fernandes, P. Glarborg, H. Hashemi, X. He, S. Mashruk, J. McGowan, C. Mounaim-Rouselle, A. Ortiz-Prado, A. Ortiz-Valera, I. Rossetti, B. Shu, M. Yehia, H. Xiao, and M. Costa,
\textit{Energy \& Fuels}, 35, (2021), 6964–7029.

\bibitem{Alnajideen24}
M. Alnajideen, H. Shi, W. Northrop, D. Emberson, S. Kane, P. Czyzewski, M. Alnaeli, S. Mashruk, K. Rouwenhorst, C. Yu, S. Eckart, and A. Valera-Medina,
\textit{Carbon Neutrality}, 3, (2024), 13.

\bibitem{El-Adawy24}
M. El-Adawy, M. A. Nemitallah, and A. Abdelhafez,
\textit{Fuel}, 364, (2024), 131090.

\end{thebibliography}{99}
\end{document}
